% ==================================================
% components_symbols.tex
% Symbol and signal-word table components
% ==================================================

% ------------------------------
% Compact symbol table for raw LaTeX renderers
% Signatures:
%   \HBSymbolTable{symbol header}{meaning header}{row macro calls}
%   \HBSymbolSignalRow{image path or basename}{optional label}{description}
%   \HBSymbolIconRow{image path or basename}{description}
% ------------------------------
\expandafter\providecommand\csname HBcomp_symbol_icon_width\endcsname{8.5mm}
\expandafter\providecommand\csname HBcomp_symbol_icon_height\endcsname{8.5mm}
% Signal pill geometry measured from the archived 1500 Ultra manual
% (MEANING OF SYMBOLS p5): all four kinds (WARNING/CAUTION/NOTE/TIP)
% share one pill of 61.2x15.4pt with the content centered both ways.
\expandafter\providecommand\csname HBcomp_symbol_signal_width\endcsname{21.5mm}
\expandafter\providecommand\csname HBcomp_symbol_signal_height\endcsname{5.4mm}
\expandafter\providecommand\csname HBcomp_symbol_signal_tri_base\endcsname{3.6}
\expandafter\providecommand\csname HBcomp_symbol_signal_tri_height\endcsname{3.2}
\expandafter\providecommand\csname HBcomp_symbol_signal_gap\endcsname{1.1mm}
\expandafter\providecommand\csname HBtype_symbol_signal_font_size\endcsname{9.9pt}
\expandafter\providecommand\csname HBtype_symbol_signal_font_leading\endcsname{10.5pt}
\providecommand{\HBSignalWordFontSize}{\csname HBtype_symbol_signal_font_size\endcsname}
\providecommand{\HBSignalWordLeading}{\csname HBtype_symbol_signal_font_leading\endcsname}
\expandafter\providecommand\csname HBcomp_symbol_icon_col_width\endcsname{14mm}
\expandafter\providecommand\csname HBcomp_symbol_signal_col_width\endcsname{26.4mm}
\expandafter\providecommand\csname HBcomp_symbol_column_gap\endcsname{2.6mm}
\expandafter\providecommand\csname HBcomp_symbol_table_tabcolsep\endcsname{2.2pt}
\expandafter\providecommand\csname HBcomp_symbol_table_row_stretch\endcsname{0.86}
\expandafter\providecommand\csname HBtype_symbol_header_font_size\endcsname{6.4pt}
\expandafter\providecommand\csname HBtype_symbol_header_font_leading\endcsname{7.0pt}
\expandafter\providecommand\csname HBtype_symbol_label_font_size\endcsname{6.4pt}
\expandafter\providecommand\csname HBtype_symbol_label_font_leading\endcsname{7.0pt}
\expandafter\providecommand\csname HBtype_symbol_body_font_size\endcsname{6.2pt}
\expandafter\providecommand\csname HBtype_symbol_body_font_leading\endcsname{7.0pt}

\providecommand{\HBTypeSymbolHeader}[1]{%
  {\HBFontSemiBold
   \color{\csname HBtype_body_color\endcsname}%
   \fontsize{\csname HBtype_symbol_header_font_size\endcsname}
            {\csname HBtype_symbol_header_font_leading\endcsname}\selectfont
   \RaggedRight #1}%
}

\providecommand{\HBTypeSymbolLabel}[1]{%
  {\HBFontSemiBold
   \color{\csname HBtype_body_color\endcsname}%
   \fontsize{\csname HBtype_symbol_label_font_size\endcsname}
            {\csname HBtype_symbol_label_font_leading\endcsname}\selectfont
   \RaggedRight #1}%
}

\providecommand{\HBTypeSymbolText}[1]{%
  {\HBFontRegular
   \color{\csname HBtype_body_color\endcsname}%
   \fontsize{\csname HBtype_symbol_body_font_size\endcsname}
            {\csname HBtype_symbol_body_font_leading\endcsname}\selectfont
   \RaggedRight #1}%
}

\newtcolorbox{hbsymboltableframe}{
  enhanced,
  clip upper,
  colback=white,
  colframe=HBSharedTableFrameColor,
  arc=\csname HBcomp_table_outer_arc\endcsname,
  boxrule=\csname HBcomp_table_outer_rule\endcsname,
  boxsep=0mm,
  left=0mm,right=0mm,top=0mm,bottom=0mm,
}

% Full-width signal-word table. The rounded frame owns only the outer stroke;
% the tabular owns the internal grid.
\newcommand{\HBSymbolTable}[3]{%
  \begingroup
  \arrayrulecolor{HBSharedTableFrameColor}
  \arrayrulewidth=\dimexpr\csname HBcomp_table_outer_rule\endcsname/2\relax
  \ifdim\arrayrulewidth<0.2pt\arrayrulewidth=0.2pt\fi
  \setlength{\tabcolsep}{\csname HBcomp_symbol_table_tabcolsep\endcsname}
  \renewcommand{\arraystretch}{\csname HBcomp_symbol_table_row_stretch\endcsname}
  \begin{hbsymboltableframe}
  \begin{tabular}{@{}|
    >{\centering\arraybackslash}m{\csname HBcomp_symbol_signal_col_width\endcsname}|
    >{\RaggedRight\arraybackslash}m{\dimexpr\linewidth
      -\csname HBcomp_symbol_signal_col_width\endcsname
      -4\tabcolsep-3\arrayrulewidth\relax}|@{}}
    \hline
    \rowcolor{HBTableHeadBg}
    \HBTypeSymbolHeader{#1} & \HBTypeSymbolHeader{#2}\tabularnewline
    \hline
    #3%
  \end{tabular}
  \end{hbsymboltableframe}
  \endgroup
  \par
}

% Independent left/right icon tables. Keeping the outer frames separate is
% what lets the V2.0 layout use six rows on the left and five on the right
% without a fake shared grid or a continuation page.
\newcommand{\HBSymbolMiniTable}[3]{%
  \begin{hbsymboltableframe}
  \begingroup
  \arrayrulecolor{HBSharedTableFrameColor}%
  \arrayrulewidth=0.35pt
  \setlength{\tabcolsep}{\csname HBcomp_symbol_table_tabcolsep\endcsname}%
  \renewcommand{\arraystretch}{\csname HBcomp_symbol_table_row_stretch\endcsname}%
  \begin{tabular}{@{}|
    >{\centering\arraybackslash}m{\csname HBcomp_symbol_icon_col_width\endcsname}|
    >{\RaggedRight\arraybackslash}m{\dimexpr\linewidth
      -\csname HBcomp_symbol_icon_col_width\endcsname
      -4\tabcolsep-3\arrayrulewidth\relax}|@{}}
    \hline
    \rowcolor{HBTableHeadBg}
    \HBTypeSymbolHeader{#1} & \HBTypeSymbolHeader{#2}\tabularnewline
    \hline
    #3%
  \end{tabular}
  \endgroup
  \end{hbsymboltableframe}%
}

\newcommand{\HBSymbolMiniTableContinuation}[1]{%
  \begin{hbsymboltableframe}
  \begingroup
  \arrayrulecolor{HBSharedTableFrameColor}%
  \arrayrulewidth=0.35pt
  \setlength{\tabcolsep}{\csname HBcomp_symbol_table_tabcolsep\endcsname}%
  \renewcommand{\arraystretch}{\csname HBcomp_symbol_table_row_stretch\endcsname}%
  \begin{tabular}{@{}|
    >{\centering\arraybackslash}m{\csname HBcomp_symbol_icon_col_width\endcsname}|
    >{\RaggedRight\arraybackslash}m{\dimexpr\linewidth
      -\csname HBcomp_symbol_icon_col_width\endcsname
      -4\tabcolsep-3\arrayrulewidth\relax}|@{}}
    \hline
    #1%
  \end{tabular}
  \endgroup
  \end{hbsymboltableframe}%
}

\newcommand{\HBSymbolTwoColumnTables}[4]{%
  \par\noindent
  \begin{minipage}[t]{\dimexpr(\linewidth-\csname HBcomp_symbol_column_gap\endcsname)/2\relax}
    \HBSymbolMiniTable{#1}{#2}{#3}%
  \end{minipage}%
  \hspace{\csname HBcomp_symbol_column_gap\endcsname}%
  \begin{minipage}[t]{\dimexpr(\linewidth-\csname HBcomp_symbol_column_gap\endcsname)/2\relax}
    \HBSymbolMiniTable{#1}{#2}{#4}%
  \end{minipage}%
  \par
}

% Localized copy can require the reference master's controlled 4+remainder
% split. Both columns preserve source order and continue on the same next page.
\newcommand{\HBSymbolTwoColumnTablesSplit}[6]{%
  \HBSymbolTwoColumnTables{#1}{#2}{#3}{#4}%
  \HBPageBreak
  \par\noindent
  \begin{minipage}[t]{\dimexpr(\linewidth-\csname HBcomp_symbol_column_gap\endcsname)/2\relax}
    \HBSymbolMiniTableContinuation{#5}%
  \end{minipage}%
  \hspace{\csname HBcomp_symbol_column_gap\endcsname}%
  \begin{minipage}[t]{\dimexpr(\linewidth-\csname HBcomp_symbol_column_gap\endcsname)/2\relax}
    \HBSymbolMiniTableContinuation{#6}%
  \end{minipage}%
  \par
}

\providecommand{\HBSymbolSignalRow}[3]{%
  \begin{minipage}[c]{\csname HBcomp_symbol_icon_col_width\endcsname}
    \centering
    \HBIfBlankTF{#2}{%
      \HBImageOrPlaceholder
        {\csname HBcomp_symbol_signal_width\endcsname}
        {\csname HBcomp_symbol_signal_height\endcsname}
        {#1}%
    }{%
      % Compact pill, content centered both ways (was a baseline-rule
      % hack that grew a tall flush slab with bottom-aligned content).
      \begingroup
      \setlength{\fboxsep}{0pt}%
      \colorbox{BrandDark}{%
        \parbox[c][\csname HBcomp_symbol_signal_height\endcsname][c]%
               {\csname HBcomp_symbol_signal_width\endcsname}{%
          \centering
          \raisebox{-0.55mm}{%
            \begin{tikzpicture}[x=1mm,y=1mm]
              \fill[white] (0,0)
                -- (\csname HBcomp_symbol_signal_tri_base\endcsname/2,\csname HBcomp_symbol_signal_tri_height\endcsname)
                -- (\csname HBcomp_symbol_signal_tri_base\endcsname,0) -- cycle;
              \node[BrandDark,font=\fontsize{4pt}{4pt}\selectfont\bfseries]
                at (\csname HBcomp_symbol_signal_tri_base\endcsname/2,1.05) {!};
            \end{tikzpicture}%
          }%
          \hspace{\csname HBcomp_symbol_signal_gap\endcsname}%
          {\color{white}\HBFontHeavy
           \fontsize{\HBSignalWordFontSize}
                    {\HBSignalWordLeading}\selectfont #2}%
        }%
      }%
      \endgroup
    }%
  \end{minipage}
  &
  \HBTypeSymbolText{#3}%
  \tabularnewline\hline
}

\providecommand{\HBSymbolIconRow}[2]{%
  \HBImageOrPlaceholder
    {\csname HBcomp_symbol_icon_width\endcsname}
    {\csname HBcomp_symbol_icon_height\endcsname}
    {#1}%
  &
  \HBTypeSymbolText{#2}%
  \tabularnewline\hline
}
